\subsection*{从输入数字到可调整的判断}
一张灰度小图可以按固定顺序展成向量，每个分量记录一个像素。要判断图中类别，最简单的尝试是给各像素分配权重，求加权和，再加一个偏置。正权重把某些像素的亮度当作支持证据，负权重当作反对证据，偏置调节整体阈值。

单个神经元接收这个分数，再通过非线性函数输出激活。以 $2x_1-x_2+1$ 为例，输入 $(1,3)$ 时分数为零，Sigmoid 输出一半。这里权重、输入、激活是三类对象：训练调整权重，样本提供输入，前向计算产生激活。

\subsection*{为什么堆很多层之前要插入非线性}
如果每层只是仿射映射，前后层可以合并：第二层展开后，仍是一个新矩阵乘原输入再加新偏置。增加深度只把同一个仿射动作拆成多步，并没有得到新的非线性边界。激活函数打断这种合并，允许不同输入区域产生不同响应。

用图形理解时，可以把第一层看成建立一组可调整的测量，后续层组合这些测量。但“某个神经元识别一条笔画”只是可能的解释，不能仅凭网络结构认定训练必然得到这种分工。

本节先把计算规则写清楚，还没有解决权重从哪里来。下一节引入损失，把“判断好不好”压缩成可优化的数，再利用微积分修改参数。理解前向计算、评价结果、更新参数三步的区别，是后续所有模型的基础。
